% ============================================================================
% A Unified Coherence Curve from Abiogenesis to 2200 CE:
% Stitching Biosphere, Civilization, and Simulated Futures
%
% Target venue: Philosophical Transactions of the Royal Society A
% Alternative: Complexity (Wiley)
% ============================================================================

\documentclass[10pt,letterpaper]{article}

\usepackage[top=0.85in,left=2.75in,footskip=0.75in]{geometry}
\usepackage[utf8]{inputenc}
\usepackage{changepage}
\usepackage{textcomp,marvosym}
\usepackage{cite}
\usepackage{nameref,hyperref}
\usepackage{amsfonts}
\usepackage{amsmath}
\usepackage{amssymb}
\usepackage{microtype}
\DisableLigatures[f]{encoding = *, family = * }
\usepackage{xcolor}
\usepackage{graphicx}
\usepackage{booktabs}
\usepackage{array}
\usepackage{float}

\usepackage[aboveskip=1pt,labelfont=bf,labelsep=period,
            justification=raggedright,singlelinecheck=off]{caption}
\usepackage[right]{lineno}
\usepackage{setspace}
\doublespacing
\pagestyle{plain}

\newcommand{\bt}{$B(t)$}
\newcommand{\kt}{$K(t)$}
\newcommand{\ct}{$C(t)$}

\begin{document}
\vspace*{0.2in}

\begin{flushleft}
{\Large\textbf{A Unified Coherence Curve from Abiogenesis to 2200 CE: Stitching Biosphere, Civilization, and Simulated Futures}}
\bigskip

Tristan Stoltz\textsuperscript{1*}

\bigskip

\textbf{1} Luminous Dynamics, Richardson, TX, USA

\bigskip

* tristan.stoltz@evolvingresonantcocreationism.com

\end{flushleft}

% ═══════════════════════════════════════════════════════════════════════════
\section*{Abstract}
% ═══════════════════════════════════════════════════════════════════════════

We present \ct{}, a unified coherence curve spanning from the origin
of life (3.8~Ga) to simulated futures (2200~CE), stitching together
three measurement regimes: biosphere coherence \bt{} (3.8~Ga to
10,000~BCE, from paleontological data), civilizational coherence \kt{}
(3000~BCE to 2020~CE, from historical proxies), and agent-based forward
projection (2026~CE to 2200~CE, from multi-world civilization simulation).

The three regimes are connected via the \textbf{headroom formula}:
\begin{equation*}
C(t) = B_{\text{HANPP}}(t) + K(t) \times (1 - B_{\text{HANPP}}(t))
\end{equation*}
which models civilization as filling the thermodynamic headroom above
the biosphere floor, rather than replacing it. \bt{} is adjusted by
Human Appropriation of Net Primary Productivity (HANPP) to reflect
the energy that civilization co-opts from the biosphere.

The unified curve reveals three insights: (1) the biosphere maintained
coherence $> 0.85$ through the pre-industrial era, and industrialization
initially depressed total coherence via HANPP before civilizational
contributions compensated; (2) the modern peak \ct{}(2020) $= 0.94$ is
the highest value in 3.8 billion years, achieved not by replacing
biological coherence but by adding civilizational organization to it;
and (3) forward projection under baseline simulation parameters shows
\ct{} stabilizing at $0.90$--$0.95$ through 2200~CE, with sensitivity
to governance quality and resource sustainability.

\noindent\textbf{Keywords:} coherence, biosphere, civilization, unified
index, HANPP, thermodynamic bridge, forward projection, deep time

\linenumbers

% ═══════════════════════════════════════════════════════════════════════════
\section*{Introduction}
% ═══════════════════════════════════════════════════════════════════════════

Measuring the organizational state of a planet---from the first
replicating molecules to interplanetary civilization---requires a
framework that spans radically different timescales, data sources, and
phenomena. Previous work has addressed pieces of this span: biosphere
coherence indices track deep-time organization from paleontological data
\cite{stoltz2026bt}; composite civilizational indices (HDI, \kt{})
measure modern and historical organization from economic and
demographic proxies \cite{stoltz2026kt, undp2020}; and simulation
models project civilizational futures under various policy assumptions
\cite{stoltz2026sim}.

No existing framework connects these three regimes into a single,
continuous, thermodynamically grounded curve. The challenge is not
merely mathematical (stitching different time series) but conceptual:
what does it mean to compare a Cambrian biosphere's organization to a
21st-century civilization's? What physical quantity bridges the
transition from biological to cultural coherence?

We propose that the bridge is \textbf{energy allocation}. Both the
biosphere and civilization compete for the same planetary energy budget
(solar insolation, chemical bonds). As civilization co-opts biosphere
energy via land use, fossil extraction, and agriculture---quantified by
Human Appropriation of Net Primary Productivity (HANPP)
\cite{haberl2007}---the biosphere's contribution to total coherence
decreases while civilization's increases. The \emph{headroom formula}
formalizes this: civilization fills the remaining organizational capacity
above the biosphere floor, producing a unified index where the two
contributions are additive rather than substitutive.

% ═══════════════════════════════════════════════════════════════════════════
\section*{Methods}
% ═══════════════════════════════════════════════════════════════════════════

\subsection*{Three Measurement Regimes}

\ct{} is computed in four phases:

\begin{enumerate}
\item \textbf{Deep time} (3.8~Ga to 10,000~BCE): Pure \bt{}, computed
from PBDB genus diversity, GEOCARB atmospheric data, $\epsilon_p$
isotopic fractionation, Bambach ecospace functional groups, and
background extinction rates. Six dimensions aggregated via geometric
mean \cite{stoltz2026bt}.

\item \textbf{Holocene transition} (10,000~BCE to 1810~CE): \bt{}
adjusted by HANPP. Before recorded K(t) data, the unified index is
simply the biosphere coherence reduced by human energy co-option:
$C(t) = B(t) \times (1 - \text{HANPP}(t))$. HANPP rises from
$\sim$0\% (pre-Neolithic) to $\sim$5\% (pre-industrial).

\item \textbf{Historical} (1810--2020~CE): The headroom formula applies.
\kt{} is computed from 35+ proxy variables across seven harmonies
\cite{stoltz2026kt}. \bt{} continues to provide the biosphere floor,
depressed by rising HANPP (5\% to 24\%).

\item \textbf{Forward projection} (2026--2200~CE): An agent-based
multi-world civilization simulator \cite{stoltz2026sim} computes
Eight Harmonies scores at each epoch. These are passed to the headroom
formula as the $K(t)$ contribution, extending \ct{} into simulated futures.
\end{enumerate}

\subsection*{The Headroom Formula}

The central equation:
\begin{equation}
C(t) = B_{\text{HANPP}}(t) + K(t) \times (1 - B_{\text{HANPP}}(t))
\end{equation}

where $B_{\text{HANPP}}(t) = B(t) \times (1 - \text{HANPP}(t))$ is the
biosphere coherence adjusted for human energy appropriation.

This formulation has three properties: (1) when $K(t) = 0$
(pre-civilizational), $C(t) = B_{\text{HANPP}}(t)$---pure biosphere;
(2) when $K(t) = 1$ (maximal civilizational coherence), $C(t) = 1$
regardless of biosphere state; (3) for intermediate values, $K(t)$
fills the ``headroom'' between the biosphere floor and perfect coherence.
Civilization cannot make \ct{} worse than the biosphere alone; it can
only add to it.

\subsection*{HANPP Trajectory}

Human Appropriation of Net Primary Productivity follows a sigmoidal
trajectory calibrated from published estimates
\cite{haberl2007, krausmann2013}:

\begin{center}
\begin{tabular}{lrr}
\toprule
Era & HANPP (\%) & $B_{\text{HANPP}}$ (if $B = 0.95$) \\
\midrule
Pre-Neolithic & 0 & 0.950 \\
Neolithic (3000~BCE) & 1 & 0.941 \\
Roman Empire (0~CE) & 3 & 0.922 \\
Medieval (1000~CE) & 4 & 0.912 \\
Industrial (1900~CE) & 12 & 0.836 \\
Modern (2020~CE) & 24 & 0.722 \\
\bottomrule
\end{tabular}
\end{center}

\subsection*{Normalization Sensitivity}

A systematic discrepancy exists between \kt{} values computed under
different normalization strategies: min-max by century yields
$K(2020) = 0.78$, while epoch-aware normalization yields $K(2020) = 0.91$
\cite{stoltz2026kt}. This propagates to \ct{}: under the first,
$C(2020) = 0.78 + 0.78 \times 0.22 = 0.95$; under the second,
$C(2020) = 0.72 + 0.91 \times 0.28 = 0.97$. We report both and note
that the qualitative shape of \ct{} is preserved across strategies.

% ═══════════════════════════════════════════════════════════════════════════
\section*{Results}
% ═══════════════════════════════════════════════════════════════════════════

\subsection*{The Unified Curve}

\ct{} spans 80 data points from 3.65~Ga to 2020~CE (Table~1).

\begin{table}[H]
\caption{{\bf Selected points on the unified coherence curve.}}
\begin{tabular}{llrrr}
\toprule
Epoch & Regime & $C(t)$ & $B_{\text{HANPP}}$ & $K(t)$ \\
\midrule
3.6~Ga (Archean) & Biosphere & 0.005 & 0.005 & --- \\
541~Ma (Cambrian) & Biosphere & 0.127 & 0.127 & --- \\
66~Ma (K-Pg) & Biosphere & 0.647 & 0.647 & --- \\
8000~BCE & Holocene & 0.888 & 0.888 & 0.00 \\
1000~CE & Holocene & 0.908 & 0.908 & 0.00 \\
1810~CE & Historical & 0.895 & 0.895 & 0.00 \\
1900~CE & Historical & 0.860 & 0.836 & 0.14 \\
1970~CE & Historical & 0.857 & 0.779 & 0.34 \\
2000~CE & Historical & 0.774 & 0.722 & 0.11 \\
2020~CE & Historical & 0.939 & 0.722 & 0.78 \\
\bottomrule
\end{tabular}
\end{table}

\subsection*{Three Insights}

\textbf{Insight 1: Industrial depression and recovery.} Between 1810
and 2000, HANPP rose from 5\% to 24\%, depressing $B_{\text{HANPP}}$
from 0.90 to 0.72. Civilizational coherence $K(t)$ rose simultaneously
from 0.0 to 0.78, but initially failed to compensate for the biosphere
loss. The nadir at 2000 ($C = 0.77$) reflects a moment when civilization
had depressed the biosphere substantially but had not yet filled the
headroom. By 2020, $K(t) = 0.78$ fills 78\% of the available headroom,
pushing $C$ back above 0.93.

\textbf{Insight 2: The 2020 peak.} $C(2020) = 0.94$ is the highest
coherence value in 3.8 billion years. This does not mean 2020 is the
``best'' era---the index measures organizational complexity, not
well-being or justice---but it means the combined biosphere-plus-civilization
system has never been more structurally integrated.

\textbf{Insight 3: The biosphere was more coherent alone.} At 1000~CE,
$C = 0.91$ with essentially zero $K(t)$ contribution. The biosphere
\emph{by itself}, before industrialization, achieved higher coherence
than most of the industrial era. This quantifies the thermodynamic cost
of the Anthropocene: civilization's rise initially \emph{reduced} total
planetary coherence before eventually exceeding it.

% ═══════════════════════════════════════════════════════════════════════════
\section*{Discussion}
% ═══════════════════════════════════════════════════════════════════════════

\subsection*{Civilization as Thermodynamic Addition}

The headroom formula's most important property is that it frames
civilization as \emph{additive}. Previous treatments of the
Anthropocene tend to frame human impact as purely destructive---biosphere
degradation, species loss, climate forcing. The \ct{} framework
acknowledges this (via HANPP depression of the biosphere floor) while
also accounting for what civilization \emph{adds}: knowledge networks,
governance structures, trade integration, reciprocal exchange. The
question ``is civilization good for the planet?'' becomes quantifiable:
\ct{} exceeds pre-civilizational $B(t)$ only when $K(t)$ fills more
headroom than HANPP destroys.

\subsection*{Connection to Consciousness}

The companion biosphere coherence paper \cite{stoltz2026bt} demonstrates
that workspace (information capacity) was the rate-limiting dimension
for consciousness emergence across 3.8~Ga. The SETI workspace filter
paper \cite{stoltz2026seti} argues this constitutes a Great Filter.
The \ct{} framework connects these findings to the present: civilizational
coherence can be understood as the \emph{continuation} of the workspace
expansion that made consciousness possible. Human civilization is the
mechanism by which the biosphere's information capacity finally exceeded
its biological ceiling---through culture, writing, computation, and
global networks rather than through genome complexity alone.

\subsection*{Limitations}

\textbf{HANPP is coarse}: Modern HANPP is estimated at 24\%
\cite{haberl2007}, but this aggregate masks regional variation
(Europe $>$ 40\%, Amazon $<$ 5\%). Regional \ct{} decomposition
is needed for policy applications.

\textbf{K(t) normalization dependence}: The 13-point divergence between
$K(2020) = 0.78$ and $K(2020) = 0.91$ means absolute \ct{} values are
normalization-dependent. Relative rankings are preserved.

\textbf{Forward projection is simulated}: The 2026--2200 phase comes
from an agent-based model, not empirical data. Its predictions are
conditional on simulation assumptions (baseline governance, continued
trade integration, no major wars).

\textbf{Single planet}: \ct{} is calibrated to Earth. Whether the
headroom formula generalizes to other planet-civilization pairings
is unknown.

% ═══════════════════════════════════════════════════════════════════════════
\section*{Conclusion}
% ═══════════════════════════════════════════════════════════════════════════

We have produced a continuous coherence curve from the origin of life
to the present, and projected it into simulated futures, using a
thermodynamically grounded bridge between biosphere and civilization.
The curve tells a story of successive organizational innovations: from
prokaryotic coherence (0.005) to multicellular integration (0.13) to
mammalian workspace expansion (0.65) to civilizational overlay (0.94).
Each transition built on---and in some cases required catastrophic
disruption of---the previous organizational regime.

The headroom formula ensures that this story is not one of replacement
but of addition. The biosphere is still there, beneath civilization,
providing the thermodynamic floor on which all human organization rests.
The 24\% of NPP that we appropriate is the price of the organizational
headroom we fill. Whether we fill it wisely is not a question the index
can answer, but it is the question the index is designed to help us ask.

% ═══════════════════════════════════════════════════════════════════════════
\begin{thebibliography}{99}
% ═══════════════════════════════════════════════════════════════════════════

\bibitem{stoltz2026bt}
Stoltz T. Catastrophe as Genesis: A thermodynamic framework for biosphere
coherence across 3.8 billion years. \emph{In review}. 2026.

\bibitem{stoltz2026kt}
Stoltz T. Historical K(t): A seven-harmony civilizational coherence index
from 3000 BCE to 2020 CE. \emph{In preparation}. 2026.

\bibitem{stoltz2026sim}
Stoltz T. Multi-world civilization simulator: 150-year agent-based
projection with consciousness-gated governance. \emph{In preparation}. 2026.

\bibitem{stoltz2026seti}
Stoltz T. The workspace filter: Why biospheres may live without
thinking. \emph{In preparation}. 2026.

\bibitem{undp2020}
UNDP. \emph{Human Development Report 2020}. United Nations; 2020.

\bibitem{haberl2007}
Haberl H, et al. Quantifying and mapping the human appropriation of net
primary production. \emph{PNAS}. 2007;104(31):12942--12947.

\bibitem{krausmann2013}
Krausmann F, et al. Global human appropriation of net primary production
doubled in the 20th century. \emph{PNAS}. 2013;110(25):10324--10329.

\bibitem{smil2017}
Smil V. \emph{Energy and Civilization: A History}. MIT Press; 2017.

\bibitem{baron2018}
Bar-On YM, Phillips R, Milo R. The biomass distribution on Earth.
\emph{PNAS}. 2018;115(25):6506--6511.

\bibitem{steffen2015}
Steffen W, et al. Planetary boundaries: guiding human development on a
changing planet. \emph{Science}. 2015;347(6223):1259855.

\end{thebibliography}

\end{document}
